\documentclass{article}
\usepackage{amsmath,amssymb}
\begin{document}

Project context: I am formalizing Higham Chapter 6, Problem 6.5 in Lean 4/mathlib.
The local repository already has complex finite matrix/SVD infrastructure for
\[
  A = U \Sigma V^*
\]
where `CMatrix n n = Fin n -> Fin n -> Complex`.

Current target:
\[
\texttt{ComplexSquareContractionMidpointProperty}(n):
\quad
\forall L,\ \|L\|_2 \le 1 \Rightarrow
\exists U,V\ \text{unitary},\quad L = (U+V)/2.
\]

Known local Lean declarations:
\begin{itemize}
\item `complexMatrixOp2 L` is the Euclidean operator norm.
\item `complexMatrixSingularValue_le_complexMatrixOp2 L i` proves every singular
  value is at most the operator norm.
\item `exists_isComplexMatrixSVD L` gives `U Sigma V` with
  `U,V : Matrix.unitaryGroup (Fin n) Complex`,
  `factorization : U * Sigma * star V = L`, and an entry formula for `Sigma`.
\item The existing `Sigma` entry formula is source-facing but, because the left
  singular-vector basis is an arbitrary extension, its nonzero rows may not be
  literally the same index as the right singular-value column. It says roughly:
  `Sigma k i = if sigma_i = 0 then 0 else if b k = leftSingularVector_i then sigma_i else 0`.
\item There is also a local theorem `complexMatrixSVDUnitary_diagonal_eq` built
  from such a left basis.
\end{itemize}

Question:
Give a rigorous finite-dimensional proof route that is friendly to Lean for
proving the midpoint theorem above. In particular:
\begin{enumerate}
\item Should I add a square-specific SVD theorem whose `Sigma` is literally
  diagonal in the same `Fin n` index, or can I use the existing arbitrary-row
  `Sigma` directly?
\item If using the existing `Sigma`, how should I define the phase matrix on
  its rows/columns so that
  \[
    \Sigma = \tfrac12(Z_+ + Z_-)
  \]
  and `Z_+`, `Z_-` are unitary, despite the row permutation/extension?
\item What are the minimal lemmas to prove in Lean first? Please spell out the
  algebraic statements, not just prose.
\item Identify any hidden caveat for zero singular values, repeated singular
  values, zero dimension, or choice of left-basis extension.
\end{enumerate}

Desired output: a proof plan with exact finite matrix identities and lemma
surfaces suitable for Lean, preferably ordered from easiest to hardest. Do not
use the spectral theorem beyond the already available SVD/singular-value facts.

\end{document}
